\section{The actual local cubic component at coefficient prime \(269\)}
\label{sec:cubic-ramified-269-component}

Put \(p=269\), \(F=K_{(269)}\), \(K=\mathbb Q(\omega)\), and
\[
 G=G_{\mathbb Q_p},\qquad H=G_F,\qquad
 \Delta=G/H=\operatorname{Gal}(F/\mathbb Q_p).
 \tag{RC.1}
\]
The extension \(F/\mathbb Q_p\) is unramified quadratic.
Thus \(I_G=I_H\), whereas the extension of profinite groups
\(1\to H\to G\to\Delta\to1\) is not assumed split.
Let \(O\) be the source coefficient ring, enlarged unramifiedly
to contain \(\mu_3\) if necessary, and \(E=\operatorname{Frac}(O)\).
All old \texttt{GD} hypotheses are retained.  In particular
\(\psi\) is finite, \(\psi^2=1\) on \(I_G\), and
\(p-1\nmid2e\).  Write
\(\gamma=\psi\kappa_p^e\), \(e\geq1\), and \(j\in\{1,2\}\).
The actual semidirect cubic seed is the restriction of AC-G49:
\[
 r_*(g).\mathrm{mat}
   =\operatorname{diag}(q(g)^j\gamma(g),1,
                         q(g)^j\gamma(g)^{-1}),\qquad
 r_*(g).\mathrm{mult}=\eta_K(g).
 \tag{RC.2}
\]
Here \(q\) is the crossed cubic character and
\(q|_H=\chi|_H\).  The semidirect involution is
\(\theta(A)=JA^{-t}J\).  The character \(\chi|_H\) is tamely
ramified of order three and satisfies \(\chi^c=\chi^{-1}\).
In particular its local factor cannot be removed by a rational
finite character, as AC-K29-269 proves.

The published local inputs are Kisin,
\emph{The structure of potentially semi-stable deformation rings},
ICM 2010 author version,
\url{https://people.math.harvard.edu/~kisin/dvifiles/licm.pdf},
\S1.1.4, equation (1.1.6), and Theorem 1.2.2,
printed pp.~5--7, and Caraiani--Newton,
\emph{On the modularity of elliptic curves over imaginary quadratic
fields}, Definition 3.3.1 and Lemma 3.3.2,
printed pp.~52--53 of
\url{https://www.ma.imperial.ac.uk/~acaraian/papers/Modularity-IQF.pdf}.
The former supplies the actual framed \(\mathrm{GL}_3\) fixed-type
quotient over \(H\), for every finite local \(E\)-algebra, including
nonreduced algebras.  Its construction is a published input;
the original construction proof is not reproduced in that survey.
The complete printed proof of the latter all-Artin ordinary lemma
is used.  Neither input is inferred from A or from the existence
of a source record bearing its name.

\begin{lemma}[AC-G65-descent: fixed-type lifting through the nonsplit quadratic quotient]
\label{AC-G65-descent}
Fix the multiplier \(\eta_K\).  Suppose the framed deformation
germ of \(r_*|_H\), with a chosen potentially crystalline Hodge
and inertial type, is formally smooth over \(E\).
Suppose this type is stable under the combined action of
conjugation of \(H\) by \(G\) and the semidirect outer involution
on matrices.  Then the germ of framed \(GG_3\) deformations of
\(r_*\) whose restriction to \(H\) has this type is formally
smooth.  Its tangent cohomology is
\[
 H^1_{\mathrm{type}}(G,M)
   \simeq H^1_{\mathrm{pcr}}(H,M)^\Delta,
 \qquad M=\operatorname{ad}r_*,
 \tag{RC.3}
\]
where the left side consists of classes whose restrictions
lie in the displayed potentially crystalline tangent subspace.
Framed tangent vectors additionally include the coboundaries.
\end{lemma}
\begin{proof}
We prove the lifting assertion, including the nonsplit issue.
Let \(B\twoheadrightarrow A\) be a small extension of finite
local \(E\)-algebras with residue \(E\), with kernel \(J_B\)
annihilated by the maximal ideal of \(B\).  In particular
\(J_B^2=0\) and \(J_B\) is an \(E\)-vector space.
Let \(r_A\) be a deformation in the asserted germ.
The obstruction to a continuous lift of its fixed-multiplier
\(GG_3\) representation lies in
\(H^2(G,M)\otimes_E J_B\).  Concretely, lift the matrix
coordinates continuously to \(B\), with the multiplier
unchanged; the defect in multiplication has values in
\(1+M_3(J_B)\).  Its first-order part is a continuous
two-cocycle for the adjoint action \(M\).  Changing the lift
changes that defect by a coboundary.  This is the usual
square-zero obstruction construction and works for the
semidirect group because its fixed component is smooth
and the outer action is defined over \(E\).

The restriction of this obstruction to \(H\) is zero:
formal smoothness of the fixed-type \(H\)-germ already
supplies an \(H\)-lift.  Restriction in degree two is
injective.  Indeed the restriction--corestriction composite
is multiplication by two, which is invertible on these
\(E\)-vector spaces.  Hence an unrestricted fixed-multiplier
\(G\)-lift \(r_0\) to \(B\) exists.

Relative to \(r_0|_H\), all \(H\)-lifts of \(r_A|_H\),
modulo conjugation by \(1+M_3(J_B)\), form the affine space
\[
                  H^1(H,M)\otimes_E J_B.
 \tag{RC.4}
\]
The classes of fixed-type lifts form a nonempty affine
subspace with translation space
\(H^1_{\mathrm{pcr}}(H,M)\otimes_E J_B\).
Here is the infinitesimal reason for this last assertion.
Use the actual framed fixed-type quotient ring.
For two lifts of a map from that ring, the difference
of every defining equation is its linear term at the
residue point: all products with \(J_B\) factor through
the residue \(E\), and products of two differences vanish.
Thus the differences of framed fixed-type lifts are exactly
the fixed-type tangent cocycles.  Strict conjugations
preserve the condition, so quotienting these differences
by the tangent coboundaries gives the claimed affine
subspace of (RC.4).

For \(\sigma\in G\setminus H\), transport an \(H\)-lift
\(\rho\) by
\[
 (T_\sigma\rho)(h)
  =r_0(\sigma)\rho(\sigma^{-1}h\sigma)r_0(\sigma)^{-1}
 \tag{RC.5}
\]
inside \(GG_3(B)\); the result has multiplier one and
is again a \(\mathrm{GL}_3\)-valued \(H\)-lift.
It fixes the origin \(r_0|_H\) of (RC.4).
Its square is conjugation of \(\rho\) by
\[
             r_0(\sigma^2)\rho(\sigma^2)^{-1}
                       \in1+M_3(J_B).
\]
Consequently (RC.5) gives an honest \(\Delta\)-action
on lift classes, although \(\sigma^2\) need not be one.
Relative to the origin it is the usual linear
\(\Delta\)-action on \(H^1(H,M)\otimes J_B\).
By the type-invariance hypothesis the nonempty fixed-type
affine subspace is stable under it.  The average of any
point of that affine subspace and its translate is
therefore a \(\Delta\)-invariant fixed-type lift class.

Finally restriction gives
\[
                   H^1(G,M)\simeq H^1(H,M)^\Delta.
 \tag{RC.6}
\]
For clarity, the inflation--restriction obstruction groups
on the two sides are cohomology groups of the finite
group \(\Delta\) with \(E\)-vector-space coefficients.
They vanish in positive degrees by averaging; this proves
both injectivity and surjectivity in (RC.6).
Lift the invariant class just obtained along (RC.6) and
modify \(r_0\) by the resulting \(G\)-cocycle.  Its
restriction is strictly conjugate to a fixed-type
\(H\)-lift, hence has the required type.  This proves
formal smoothness.  The same restriction argument on
dual numbers gives (RC.3).
\end{proof}

\begin{theorem}[AC-G65-local: the actual twelve-dimensional component at \(269\)]
\label{AC-G65-local}
Let \(U\) be the unrestricted framed local \(O\)-algebra
for the reduction of (RC.2), with multiplier fixed to
\(\eta_K\), and \(x_*:U\to O\) its actual basepoint.
There is a prime \(P_*\subset\ker x_*\) and an injective
formal chart
\[
 \Phi:U/P_*\hookrightarrow E[[X_1,\ldots,X_{12}]]
 \tag{RC.7}
\]
with constant term \(x_*\), having weights
\(-2,-1,0,1,2\) with multiplicities
\[
                         (1,2,3,4,2).
 \tag{RC.8}
\]
All actual matrix entries have their usual matrix weights,
and their linear coefficients span the coordinates.
The chart has the saturated cotangent property: if
\(\Phi(f)\in(X)^n\), then
\[
 \exists\,0\ne a\in O:\quad
       af\in(\ker x_*)^n+P_*.
 \tag{RC.9}
\]
Killing \(P_*\) is preserved by every strict conjugation
over every complete Noetherian local source \(O\)-test
algebra, including algebras with \(p\)-torsion or
nilpotents.  Every finite local \(E\)-algebra point
killing \(P_*\) restricts to a potentially crystalline
representation of \(H\) of the type of \(r_*|_H\).
Conversely every such point specializing to \(x_*\)
kills \(P_*\).
\end{theorem}
\begin{proof}
\proofstep{The actual restriction quotient and its coefficient domains.}
The arbitrary-residual framed construction of AC-G49
gives \(U\); it does not require the residual matrix
to come from one rational character.  Restrict its
universal representation to \(H\).  There its multiplier
is one, so the framed \(\mathrm{GL}_3\) universal property
gives a map from the corresponding framed \(H\)-ring to
\(U\), over the same \(O\).
Take the potentially crystalline inertial and Hodge
types of
\[
 r_*|_H=\operatorname{diag}(\chi^j\gamma,1,
                                      \chi^j\gamma^{-1}).
 \tag{RC.10}
\]
The finite type has characters
\(\chi^j\psi,1,\chi^j\psi^{-1}\), and the Hodge type is
parallel with the weights of
\(\kappa_p^e,1,\kappa_p^{-e}\) at both embeddings of \(F\).
These data are defined over \(E\).  The parallel Hodge
type has rational reflex field; enlarging \(E\) for
auxiliary embeddings in a comparison argument does
not change this defining type.

Pull the Kisin fixed-type ideal to \(U\), and let
\(I_{\mathrm{cr}}\) be its \(p\)-saturation.
For every finite local \(E\)-algebra \(C\), an
\(O\)-algebra map \(U\to C\) kills \(I_{\mathrm{cr}}\)
if and only if its restriction to \(H\) has the
specified potentially crystalline type: saturation
has no effect because \(p\) is a unit in \(C\).
The integral basepoint kills \(I_{\mathrm{cr}}\)
because \(O\) is \(p\)-torsion free.
This gives an actual closed quotient, rather than
a condition imposed only on field-valued points.

All finite Artin points of its formal germ are
continuous.  As in AC-G54, present \(U\) by finitely
many power-series variables after subtracting their
basepoint values.  At a finite Artin point at \(x_*\)
these variables are nilpotent; evaluation factors
through a fixed finite power of their ideal and is
continuous in each of its finitely many coefficient
maps.  The same observation applies after adjoining
formal conjugation variables.

\proofstep{Smoothness over \(H\) and invariance of the type.}
On \(H\) the three diagonal characters of (RC.10)
are pairwise distinct: a nonzero cyclotomic power
has infinite image on every open inertia subgroup.
Thus the adjoint invariants have dimension three.
The de Rham tangent quotient has dimension six,
three positive root directions at each of the two
embeddings of \(F\).
Kisin's complete two-term complex in \S1.1.4 and
equation (1.1.6) therefore give
\[
 \dim H^1_{\mathrm{pcr}}(H,M)=6+3=9.
\]
Adding the six framed coboundaries gives tangent
dimension fifteen.  Theorem 1.2.2 gives a
fifteen-dimensional formally smooth component
through this point.  Hence the local ring has both
dimension and embedding dimension fifteen and
is regular; its completion is \(E[[T_1,\ldots,T_{15}]]\).
This checks smoothness of the full germ, not just
of one component while ignoring intersections.

For an element outside \(H\), field conjugation
replaces \(\chi\) by \(\chi^{-1}\) and leaves the
rational characters \(\psi,\kappa_p\) unchanged.
The matrix involution sends
\[
 \operatorname{diag}(\chi^{-j}\psi,1,\chi^{-j}\psi^{-1})
       \ \longmapsto\
 \operatorname{diag}(\chi^j\psi,1,\chi^j\psi^{-1}).
 \tag{RC.11}
\]
Conjugation by the remaining diagonal basepoint
matrix does not change the type.  The same
inverse-and-reverse operation restores the ordered
Hodge powers \((e,0,-e)\), while the two embeddings
are interchanged.  Since the Hodge type is parallel,
both types are invariant under exactly the combined
action in AC-G65-descent.  That lemma now proves
formal smoothness of the \(GG_3\) germ of
\((U/I_{\mathrm{cr}})[1/p]\) at \(x_*\).

\proofstep{The actual tangent dimensions and grading.}
The weight torus is \(T(t)=\operatorname{diag}(t,1,t^{-1})\).
It satisfies \(\theta(T)=T\) and centralizes (RC.2)
in the semidirect sense.  Thus its weights commute
with the \(G\)-action on \(M\).
As \(G\)-modules the weights are
\[
\begin{array}{c|ccccc}
i&-2&-1&0&1&2\\ \hline
M_i&
\eta_K\gamma^{-2}&
\gamma^{-1}\operatorname{Ind}_H^G\chi^j&
1\oplus\eta_K\oplus\eta_K&
\gamma\operatorname{Ind}_H^G\chi^j&
\eta_K\gamma^2 .
\end{array}
 \tag{RC.12}
\]
An additional \(\eta_K\) on either induced summand
gives an isomorphic module, by changing the sign
of one of its two basis vectors.  Formula (RC.12)
also follows directly from
\(X\mapsto-r_*(g).\mathrm{mat}\,JX^tJ\,
             r_*(g).\mathrm{mat}^{-1}\) outside \(H\).
The only adjoint invariants are the weight-zero
line \(\operatorname{diag}(a,0,-a)\).

The two-term potentially crystalline complex over
\(H\) is equivariant for \(\Delta\).  Taking
\(\Delta\)-invariants is exact, since two is invertible.
In each weight its two copies of the period module
cancel in the dimension calculation, so (RC.3) gives
\[
 \dim H^1_{\mathrm{type}}(G,M_i)
   =\dim\bigl(D_{\mathrm{dR},H}(M_i)/
                 \operatorname{Fil}^0\bigr)^\Delta
                      +\dim M_i^G.
 \tag{RC.13}
\]
The first term is \(0,0,0,2,1\) in weights
\(-2,-1,0,1,2\).  Indeed the positive cyclotomic
weights are \(e,e,2e\); in the de Rham module over
the unramified quadratic field the two embeddings
are exchanged, and descent leaves one copy of
each rational positive-weight direction.
The second term is \(0,0,1,0,0\).
The coboundary dimensions are respectively
\(1,2,2,2,1\), obtained by subtracting the invariant
line from the nine-dimensional adjoint space.
Their sums are exactly (RC.8).  Thus the formally
smooth completed \(GG_3\) germ is \(E[[X_1,\ldots,X_{12}]]\).

\proofstep{Select the actual prime and pass back to all integral tests.}
Let \(P_*\) be the kernel of the map from \(U\) to
this completed local ring.  It is prime and
\(I_{\mathrm{cr}}\subset P_*\subset\ker x_*\).
Denominator clearing at the \(E\)-valued point,
as fully proved in AC-G54, gives (RC.9), and gives
a single \(y\in U\) with
\[
           x_*(y)\ne0,\qquad yP_*\subset I_{\mathrm{cr}}.
 \tag{RC.14}
\]
The latter follows because the finitely generated
module \(P_*/I_{\mathrm{cr}}\) vanishes after
localization at the regular germ.  Equations
(RC.14) and the nonzero residue of \(y\) prove the
two assertions about finite Artin points.

For completeness the integral conjugation argument
uses a universal test, not an embedding of an
arbitrary torsion ring into characteristic zero.
Put \(R=U/P_*\), adjoin nine variables \(Z_{ab}\),
and conjugate the universal representation by
\(1+(Z_{ab})\) over \(R[[Z]]\).
The unrestricted framed property gives
\(a:U\to R[[Z]]\).
Embed this target coefficientwise into \(E[[X,Z]]\).
Modulo every power of \((X,Z)\), its representation
is conjugate to a fixed-type Artin point, so
Kisin's full Artin assertion gives \(a(I_{\mathrm{cr}})=0\).
The value \(a(y)\) is a unit in \(E[[X,Z]]\);
(RC.14) then gives \(a(P_*)=0\).
The coefficientwise embedding is injective, hence
this equality already holds in \(R[[Z]]\).
Evaluation at the entries of any strict conjugating
matrix over any source complete local \(O\)-algebra
proves the assertion there, including \(p\)-torsion
and nilpotents.

Finally choose finitely many actual entry differences
whose differentials form bases of the five tangent
weight spaces.  They are homogeneous because of
the exact weight-torus action just described.
The formal inverse-function theorem makes them
coordinates of the regular completed germ.
Equivariance then makes every actual entry
homogeneous of its matrix weight; a nonzero
basepoint entry has weight zero.
This proves spanning, grading and the complete
saturated chart assertions.
\end{proof}

\begin{theorem}[AC-G65-flag: the exact integral inertia flag at \(269\)]
\label{AC-G65-flag}
Let \(R=U/P_*\).  For every commutative local
\(O\)-algebra \(A\) with the prescribed residue map,
and every specialization \(r_A\) of its universal
representation, there is \(h\in\mathrm{GL}_3(A)\),
\(h\equiv1\pmod{\mathfrak m_A}\), such that
\(h^{-1}r_Ah\) is upper triangular on \(H\).
Its ordered inertia characters are exactly
\[
                 \chi^j\gamma,\qquad1,\qquad
                                  \chi^j\gamma^{-1}.
 \tag{RC.15}
\]
Thus the flag is integral and holds also over
characteristic-\(269\) complete DVRs.  The cubic
factor in (RC.15) is part of the assertion; this
is not the old inertia condition with that factor
silently omitted.
\end{theorem}
\begin{proof}
\proofstep{All Artin points at the actual germ.}
Let \(C\) be a finite local \(E\)-algebra with residue
\(E\), and \(r_C\) an Artin point at \(x_*\).
Its restriction to \(H\) has the fixed potentially
crystalline type by AC-G65-local.  Choose a finite
Galois extension \(L/F\) over which this restriction
is crystalline and which kills both finite
characters \(\chi\) and \(\psi\).
After a harmless finite coefficient extension to
include the embeddings of \(L\), Caraiani--Newton
Lemma 3.3.2 applies to \(G_L\).
The residue is the direct sum of the three
ordinary characters \(\kappa_p^e,1,\kappa_p^{-e}\).
The Hodge type is parallel and regular; in the
convention of their Definition 3.3.1 its dominant
weight is \((e-2,-1,-e)\), which is dominant for
every \(e\geq1\).
Their all-Artin conclusion supplies an ordered
ordinary flag with those inertia characters.

Choose \(s\in I_L\) for which
\(\kappa_p(s)^e,1,\kappa_p(s)^{-e}\) are distinct.
The cyclotomic image on every open inertia subgroup
is infinite, so such an element exists.
The Lagrange projectors of \(r_C(s)\) give three
rank-one summands, and any flag with the specified
ordered characters is the ordered sum of these
summands.  The projectors have coefficients in
\(C\), so they descend the flag from the auxiliary
coefficient extension and prove uniqueness.
Since \(L/F\) is Galois, conjugation by \(H\)
preserves the prescribed ordered inertia
characters on \(I_L\).  Uniqueness therefore
makes the flag \(H\)-stable.

Divide the ordered inertia quotient characters
by their respective powers of \(\kappa_p\).
They are trivial on \(I_L\), hence have finite
image.  Their reductions are respectively
\(\chi^j\psi,1,\chi^j\psi^{-1}\), from the seed.
After division by those constant characters
their values lie in \(1+\mathfrak m_C\).
The finite nilpotent logarithm identifies this
group with the torsion-free additive
\(E\)-vector space \(\mathfrak m_C\).
Their finite images are consequently trivial.
This proves exactly (RC.15) for every such Artin
point, including its finite characters.

\proofstep{Polynomial identities in the actual integral ring.}
Put \((a_0,a_1,a_2)=(\chi^j\gamma,1,\chi^j\gamma^{-1})\)
on \(H\).  Apply the preceding conclusion to
the formal chart modulo every power of its
maximal ideal.  Triangular matrices then give
\[
\begin{aligned}
 \det(X-r(s))&=\prod_{i=0}^2(X-a_i(s))
                                      &&(s\in I_G),\\
 \prod_{i=0}^2\bigl(r(s_i)-a_i(s_i)I\bigr)&=0
                                      &&(s_i\in I_G),
\end{aligned}
 \tag{RC.16}
\]
where the factors in the second line are in
increasing order from left to right.
The rightmost factor lowers the flag one step,
then the middle factor does so, and the leftmost
kills the first step.  The formal series ring
is separated and the chart is injective, so
(RC.16) holds in \(R\) itself and survives every
\(O\)-algebra specialization.  This uses all
Artin points of the germ, not a density claim
about unrelated classical points.

\proofstep{A residual unit witness despite ramification of \(\chi\).}
Choose \(s_0\in I_G\) with
\(\overline\kappa_p(s_0)^{2e}\ne1\), using
\(p-1\nmid2e\).
On inertia \(\psi^2=1\), so
\(\overline\gamma(s_0)^2\ne1\).
The value \(\overline\gamma(s_0)\) has order
dividing \(268\): the finite inertia part of
\(\psi\) is quadratic and the mod-\(p\)
cyclotomic values lie in \(\mathbb F_p^\times\).
Because \(\gcd(3,268)=1\), the element \(s=s_0^3\)
satisfies
\[
 \chi(s)=1,\qquad
              \gamma(s)^2-1\in O^\times.
 \tag{RC.17}
\]
The three seed values at \(s\) are consequently
\(d_0=\gamma(s),d_1=1,d_2=\gamma(s)^{-1}\),
pairwise distinct modulo the maximal ideal of
every allowed \(A\).

\proofstep{The strict flag over arbitrary local rings.}
Let \(M=r_A(s)\).  Its Lagrange projectors
\[
 E_i=\prod_{k\ne i}\frac{M-d_k I}{d_i-d_k}
 \tag{RC.18}
\]
are defined over \(A\), are orthogonal, and reduce
to the coordinate projectors.  Thus
\(h=(E_0e_0,E_1e_1,E_2e_2)\equiv I\) is invertible
and diagonalizes \(M\) to \(D=\operatorname{diag}(d_i)\).
For \(x\in I_G\), write \(N=r_A(x)\) in that basis.
The ordered identity (RC.16), at
\((x,s,s),(s,x,s),(s,s,x)\), gives
\[
\begin{split}
 (N-a_0(x)I)(D-d_1I)(D-d_2I)&=0,\\
 (D-d_0I)(N-a_1(x)I)(D-d_2I)&=0,\\
 (D-d_0I)(D-d_1I)(N-a_2(x)I)&=0.
\end{split}
 \tag{RC.19}
\]
The first identity determines the first column
below the diagonal and its diagonal entry;
the third determines the last row and its
diagonal entry; the \((1,1)\) entry of the
second gives the middle diagonal entry.
All cancelled differences \(d_i-d_k\) are units.
Thus all inertia matrices are upper triangular
with diagonal (RC.15).

For \(g\in H\), the characters \(a_i\) are
characters of \(H\), so \(r_A(gsg^{-1})\) is
upper triangular with diagonal \(d_i\).
The identity
\(r_A(gsg^{-1})r_A(g)=r_A(g)D\)
successively forces entries \((2,0),(2,1),(1,0)\)
of \(r_A(g)\) to vanish, again by unit
differences.  This is upper triangularity on
the exact subgroup \(H=\ker\eta_K\).
No domain, reducedness or characteristic-zero
assumption on \(A\) was used in this step.
\end{proof}

\begin{corollary}[AC-G65-model: the actual model and strict integral extraction at \(269\)]
\label{AC-G65-model}
Retain the same quantified source inputs as AC-G49--G52,
in particular \texttt{CohInp}, \texttt{LemmaP}, and
\texttt{LemmaM'} where those results use them.
Use \(P_*\) of AC-G65-local as the actual local condition
at \(269\).  Then the arbitrary-residual global
construction, the relative relation count, the complete
graded-model properties and the strict integral
extraction of AC-G49--G52 hold also at coefficient
prime \(269\).
In particular, for this actual model \(\mathcal M_j\),
\[
\begin{gathered}
 \alpha_j=\chi^{-j}\psi^{-1}\kappa_{269}^{-e}|_{G_K},
 \qquad \operatorname{length}_O\mathcal S_j
                    \leq\mathcal M_j.\mathrm{len},\\
 \mathcal S_j=\ker\!\left(
 \begin{aligned}
 &H^1(G_K,(E/O)(\alpha_j))\longrightarrow\\
 &\quad\prod_{w\mid269}H^1(G_{K_w},(E/O)(\alpha_j))
 \times\prod_{w\nmid269}H^1(I_w,(E/O)(\alpha_j))
 \end{aligned}\right).
\end{gathered}
 \tag{RC.20}
\]
The same bound holds after imposing any additional
decomposition-group vanishing conditions.
This is a comparison with the newly constructed
model, not a numerical upper bound for its length.
\end{corollary}
\begin{proof}
We check the places where the earlier stated prime
range was used.
The global cubic basepoint remains integral and
normalized, and its reduction detects \(\mu_3\)
because \(269\ne3\).
At every prime away from the coefficient prime,
the kernel of its residual inertia still kills
the characteristic-zero basepoint: the finite
orders involved divide
\(2\cdot3\cdot\operatorname{ord}(\psi)\), prime
to \(269\), while the cyclotomic character is
unramified there.  The arbitrary framed
\(GG_3\) construction and actual quotient by
these inertia conditions in AC-G49 therefore
apply, using \(P_*\) as their local ideal.
Their only local requirements are an actual
universal local ring, the basepoint on the
chosen quotient, and strict conjugation
stability; AC-G65-local supplies all of them.

For AC-G50, the away-from-\(269\) residual
inertia indices are prime to \(269\), so its
literal \texttt{IsLocCond} proof and generic
cohomological count are unchanged.
Its global cyclotomic witness remains valid
after cubing, since cubing kills the cubic
character and preserves a cyclotomic value
\(-1\).  Its correction at infinity uses only
\(269\ne2\).
The local chart has exactly the five
multiplicities used in AC-G50--G51, by (RC.8).
The entry grading, eight normalization
equations, diagonal integral derivative
calculation and saturated local cotangent
property therefore supply every corresponding
field of the actual graded model.  These
arguments quantify over the new actual
basepoint and local ring; no old scalar
transport is asserted.

The only additional local step in AC-G52
used unramifiedness of \(\chi\) to choose
an inertia element with \(\chi(s)=1\) and
\(\gamma(s)^2-1\) a unit.
Equation (RC.17) proves precisely those
two requirements at \(269\).
Consequently its local cobounding vector
is again killed by the same torsion scalar:
on the weight-\(-1\) induced lattice the
action of \(s\) is \(\gamma(s)^{-1}I\),
and \(1-\gamma(s)^{-1}\) is a unit.
The normalized integral cocycle, actual
square-zero global point and compatible
local conjugation lift are exactly the
ones proved in AC-G52, now using the
strict conjugation stability of \(P_*\).
Its injective family of derivations and
the expressly retained \texttt{LemmaM'}
give (RC.20).  The passage from bounded
torsion to the full strict group and the
extra-local-condition inclusion are the
same exact directed-union argument there.
\end{proof}

\paragraph{Point transfer and remaining arithmetic scope.}
For an actual Type II point with the middle-character
and endpoint-determinant identities retained in
AC-K29-point, the endpoint block over \(H\) has
ordered inertia \(\chi^j\gamma,\chi^j\gamma^{-1}\).
Untwisting that block by \(\chi^{-j}\) removes
the cubic inertia and gives \(\gamma,\gamma^{-1}\).
Its rational extension from AC-G45 therefore has
an ordinary flag on \(G\): the flag on \(H\) is
unique on open inertia and its ordered characters
are rational characters, so conjugation by \(G\)
preserves it.  Its potentially crystalline
inertial type is now \(\psi\oplus\psi^{-1}\).
On the specific source subgroup
\(H_L=\ker\eta_K\cap\ker\tau\), with
\(\psi_I=\tau_I\), this type is trivial.
Finite Galois descent of its full crystalline
period module, exactly as in the complete
AC-G57-TypeII argument, makes it crystalline
on that subgroup.  Thus no cubic inertia
factor is incorrectly retained in the old
endpoint condition, or incorrectly deleted
from the full rank-three condition.

The integral local model and strict extraction
at this previously omitted prime are now supplied.
They do not prove the actual global finite-support
condition, Type I automorphy, mixed tangent
vanishing, or a numerical bound for
\(\mathcal M_j.\mathrm{len}\).
In particular the ramified CM factor at \(269\)
still has \(D_{\mathrm{cris}}=0\) and \(U_{269}=0\),
as AC-K29-269 proves.  The ordinary and nonzero
finite-slope Euler constructions used at the
other coefficient primes do not follow from
this new local deformation component.

\paragraph{Geometric dictionary.}
The construction is a fixed-multiplier deformation
problem for a rank-three local system, with a
closed potentially crystalline condition imposed
after the unramified quadratic cover and then
descended through its order-two automorphism
group.  The strict integral flag consists of
free direct summand subbundles defined by the
explicit polynomial projectors (RC.18).
The already audited SGA~1, Expos\'e~V,
\S\S4 and~7, supplies only the finite-cover and
finite-local-system categorical dictionary.
The characteristic-zero fixed-type theorem,
the nonsplit infinitesimal descent, and the
integral model-length comparison have the
separate proofs and published inputs above.
